% from ega4-18-fr.tex
\oldpage[IV-4]{137}

on a vu qu'un $A$-homomorphisme de $A_1$ dans $A_2$ est nécessairement local; il correspond à un $S$-morphisme $f$ de $X'_2=\Spec(A_2)$ dans $X_1$ tel que $f(x_2)=x_1$, ou encore, en posant $X'_3=X_1\times_S X'_2$, à une $X'_2$-section $f'$ de $X'_3$ telle que $f'(x_2)=x_3$. Comme $k(x_3)=k(x_2)$ et que $X'_2$ est connexe, l'unicité de $f'$ résulte de (17.4.9), d'où (i).

\begin{env}[18.6.4]
\label{IV.18.6.4}
Désignons par $\mathfrak S$ le sous-ensemble de l'ensemble $\mathfrak E$ défini dans (18.6.1) formé des $A$-algèbres strictement essentiellement étales appartenant à $\mathfrak E$. Il résulte de (18.6.3) que la relation « il existe un $A$-homomorphisme de $A_1$ dans $A_2$ » est une relation de \emph{préordre} dans $\mathfrak S$ et fait de $\mathfrak S$ un ensemble \emph{filtrant croissant}. En indexant $\mathfrak S$ par lui-même au moyen de l'application identique, il résulte de (18.6.3, (i)) que si $\lambda\leq\mu$ dans $\mathfrak S$, il existe un $A$-homomorphisme \emph{unique} $\varphi_{\mu\lambda}:A_\lambda\to A_\mu$ et $(A_\lambda,\varphi_{\mu\lambda})$ est évidemment un \emph{système inductif} de $A$-algèbres locales, les $\varphi_{\mu\lambda}$ étant des homomorphismes locaux. En outre, il résulte de (17.3.5) que pour $\lambda\leq\mu$, $A_\mu$ est une $A_\lambda$-algèbre strictement essentiellement étale.
\end{env}

\begin{definition}[18.6.5]
\EGAReaderTarget{ega4.c04.def.18-6-5}{IV.18.6.5}
On appelle \emph{hensélisé} d'un anneau local $A$ et on désigne par ${}^h\!A$ la $A$-algèbre limite inductive du système inductif $(A_\lambda,\varphi_{\mu\lambda})$ défini dans (18.6.4).
\end{definition}

Cette définition ne dépend qu'en apparence du choix de $\mathfrak E$; si $\mathfrak E'$ est un autre ensemble de $A$-algèbres locales essentiellement étales ayant la même propriété que $\mathfrak E$, $\mathfrak S'$ le sous-ensemble de $\mathfrak E'$ formé des $A$-algèbres strictement essentiellement étales, il résulte de (18.6.3, (ii)) que, pour les relations de préordre envisagées, $\mathfrak S$ et $\mathfrak S'$ sont \emph{cofinaux} à $\mathfrak S''=\mathfrak S\cup\mathfrak S'$, donc donnent la même limite inductive à un isomorphisme près. Nous allons voir d'ailleurs ci-dessous (18.6.6, (i) et (ii)) que ${}^h\!A$ et l'homomorphisme canonique $A\to{}^h\!A$ forment une solution d'un problème universel, et par suite sont déterminés à isomorphisme unique près.

On remarquera que si $A$ est un corps, on a évidemment ${}^h\!A=A$. Dans le cas général, ${}^h\!A$ est aussi le hensélisé de toutes les $A$-algèbres strictement essentiellement étales $A_\lambda$ ($\lambda\in\mathfrak S$), car si $B$ est une $A_\lambda$-algèbre strictement essentiellement étale, $B$ est aussi une $A$-algèbre strictement essentiellement étale (18.6.1), donc (à un $A$-isomorphisme près) une des $A_\mu$ pour $\mu\geq\lambda$, et les $A_\mu$ telles que $\mu\geq\lambda$ et que $A_\mu$ soit une $A_\lambda$-algèbre strictement essentiellement étale forment un ensemble cofinal dans l'ensemble préordonné des $A_\lambda$, en vertu de (18.6.1) et (18.6.3).

\begin{theorem}[18.6.6]
\label{IV.18.6.6}
Soient $A$ un anneau local, ${}^h\!A$ son hensélisé.
\begin{enumerate}
  \item[(i)] ${}^h\!A$ est un anneau local hensélien et l'homomorphisme structural $A\to{}^h\!A$ est local.
  \item[(ii)] Pour tout anneau local hensélien $B$, l'application canonique
  \[
    \Hom.\operatorname{loc}({}^h\!A,B)\longrightarrow\Hom.\operatorname{loc}(A,B)
  \]
  est bijective.
  \item[(iii)] ${}^h\!A$ est un $A$-module fidèlement plat, et si $\mathfrak m$ est l'idéal maximal de $A$, $\mathfrak m.{}^h\!A$ est l'idéal maximal de ${}^h\!A$, et l'homomorphisme $A/\mathfrak mA\to{}^h\!A/\mathfrak m.{}^h\!A$ des corps résiduels est bijective.
  \item[(iv)] Si $\widehat A$ et $({}^h\!A)\widehat{\phantom{x}}$ sont les séparés complétés des anneaux locaux $A$ et ${}^h\!A$, l'homomorphisme $\widehat A\to({}^h\!A)\widehat{\phantom{x}}$ déduit de l'homomorphisme structural $A\to{}^h\!A$ par complétion est bijectif.
  \item[(v)] Pour que ${}^h\!A$ soit noethérien, il faut et il suffit que $A$ le soit.
  \item[(vi)] Si $A$ est hensélien, l'homomorphisme canonique $A\to{}^h\!A$ est bijectif.
\end{enumerate}
\end{theorem}

